%==============================================================================
% 问题描述与动力学建模
%==============================================================================
\section{问题描述与动力学建模}\label{sec:model}

%------------------------------------------------------------------------------
% 坐标系定义
%------------------------------------------------------------------------------
\subsection{坐标系与状态定义}

本文采用地面惯性坐标系作为参考坐标系，飞行器运动用三自由度（3-DOF）质点模型描述。该模型忽略飞行器姿态动力学，将飞行器视为受控质点，适用于轨迹层面的控制分析与设计。

定义状态向量 $\bx \in \Real^6$ 和控制向量 $\bu \in \Real^3$ 如下：
\begin{equation}\label{eq:state_control}
\bx = \begin{bmatrix} p_x & p_y & p_z & V & \chi & \gamma \end{bmatrix}\trans, \quad
\bu = \begin{bmatrix} T & \dot{\chi} & \dot{\gamma} \end{bmatrix}\trans
\end{equation}
其中，$p_x, p_y, p_z$ 为飞行器在惯性系中的位置坐标（m）；$V$ 为飞行速度大小（m/s）；$\chi$ 为航向角（rad），定义为速度向量在水平面内的投影与$x$轴正方向的夹角；$\gamma$ 为航迹角（rad），定义为速度向量与水平面的夹角。控制量中，$T$ 为发动机推力（N），$\dot{\chi}$ 为航向角速率（rad/s），$\dot{\gamma}$ 为航迹角速率（rad/s）。

%------------------------------------------------------------------------------
% 动力学方程
%------------------------------------------------------------------------------
\subsection{三自由度质点动力学方程}

在地面惯性坐标系下，飞行器的三自由度质点运动方程为：
\begin{equation}\label{eq:dynamics}
\dot{\bx} = f(\bx, \bu) = 
\begin{bmatrix}
V \cos\gamma \cos\chi \\[4pt]
V \cos\gamma \sin\chi \\[4pt]
V \sin\gamma \\[4pt]
\dfrac{T - D(\bx)}{m} - g \sin\gamma \\[6pt]
\dot{\chi} \\[4pt]
\dot{\gamma}
\end{bmatrix}
\end{equation}
其中，$m$ 为飞行器质量（kg）；$g$ 为重力加速度（m/s$^2$）；$D(\bx)$ 为气动阻力，采用如下二次型模型：
\begin{equation}\label{eq:drag}
D(\bx) = \frac{1}{2} \rho V^2 S C_D
\end{equation}
式中，$\rho$ 为空气密度（kg/m$^3$），$S$ 为参考面积（m$^2$），$C_D$ 为阻力系数。

式~\eqref{eq:dynamics}的前三个分量描述位置运动学，第四个分量描述速度动力学（含推力、阻力与重力分量），最后两个分量将航向角速率和航迹角速率作为直接控制输入。该模型的特点是：位置运动学具有三角函数非线性，速度动力学含有重力耦合项，而角速率通道为积分关系。

%------------------------------------------------------------------------------
% 负推力建模
%------------------------------------------------------------------------------
\subsection{负推力建模}\label{sec:neg_thrust}

在实际飞行中，飞行器减速除依靠减小发动机推力外，还可通过展开减速板（speedbrake）产生附加阻力实现快速减速。为在控制模型中反映这一能力，本文将减速板的减速效果等效为负推力，从而扩展推力的下限。

设发动机自身推力范围为 $[T_{\text{eng,min}},\, T_{\text{eng,max}}]$，减速板可产生的等效负推力为 $T_{\text{brake}}$，则等效推力的约束为：
\begin{equation}\label{eq:thrust_constraint}
T_{\min} \leq T \leq T_{\max}
\end{equation}
其中 $T_{\min} = T_{\text{eng,min}} - T_{\text{brake}}$，$T_{\max} = T_{\text{eng,max}}$。本文取 $T_{\min} = -50$\,N，$T_{\max} = 500$\,N，负推力下限对应减速板全开状态。

%------------------------------------------------------------------------------
% 控制约束
%------------------------------------------------------------------------------
\subsection{约束条件}

飞行器运动受以下约束：
\begin{equation}\label{eq:constraints}
\begin{cases}
T_{\min} \leq T \leq T_{\max} \\[4pt]
|\dot{\chi}| \leq \dot{\chi}_{\max} \\[4pt]
|\dot{\gamma}| \leq \dot{\gamma}_{\max} \\[4pt]
V_{\min} \leq V \leq V_{\max}
\end{cases}
\end{equation}
上述约束包括推力幅值约束、航向角速率与航迹角速率约束以及速度约束。这些约束在MPC优化问题中以不等式形式显式处理。

%------------------------------------------------------------------------------
% 轨迹跟踪问题描述
%------------------------------------------------------------------------------
\subsection{轨迹跟踪问题}

轨迹跟踪的目标是设计控制律 $\bu(t) = \pi(\bx(t),\, \bx_r(t))$，使得飞行器状态 $\bx(t)$ 在存在初始偏差的情况下渐近跟踪参考轨迹 $\bx_r(t)$，同时满足约束条件~\eqref{eq:constraints}。形式化地，定义跟踪误差 $\tilde{\bx}(t) = \bx(t) - \bx_r(t)$，目标是使
\begin{equation}
\lim_{t \to \infty} \|\tilde{\bx}(t)\| \to 0
\end{equation}
且控制输入始终满足 $\bu(t) \in \mathcal{U}$，其中 $\mathcal{U}$ 为由式~\eqref{eq:constraints}定义的可行控制集。

本文采用LTV-MPC方法求解该问题，其核心思路是在参考轨迹附近对非线性动力学~\eqref{eq:dynamics}进行线性化，将非线性的有限时域最优控制问题转化为每步可快速求解的QP问题。具体设计见第~\ref{sec:mpc}节。
